\documentclass[../cours_spe_trigo.tex]{subfiles}
\begin{document}
\setcounter{chapter}{6}
\chapter{Les formules d'addition, démontrées avec les rotations}
\section{Comment une rotation agit-elle sur les coordonnées ?}
Une rotation du plan conserve les longueurs et les angles orientés. La rotation d'angle $a$ envoie le vecteur $\vec{OI}=(1,0)$ sur $(\cos a,\sin a)$. Elle envoie le vecteur $\vec{OJ}=(0,1)$, qui lui est perpendiculaire et orienté d'un quart de tour positif, sur $(-\sin a,\cos a)$.
\begin{demonstration}
Tout vecteur $(x,y)$ s'écrit $x(1,0)+y(0,1)$. Une rotation étant linéaire, son image est
\[
\begin{aligned}
R_a(x,y)&=x(\cos a,\sin a)+y(-\sin a,\cos a)\\
&=(x\cos a-y\sin a,\ x\sin a+y\cos a).
\end{aligned}
\]
Ainsi la matrice de rotation est
\[\boxed{R_a=\begin{pmatrix}\cos a&-\sin a\\\sin a&\cos a\end{pmatrix}.}\]
Cette formule est obtenue à partir de la géométrie des deux vecteurs de base, \emph{sans employer les formules d'addition que nous allons prouver}. $\square$
\end{demonstration}
\section{Composition de rotations et addition des angles}
Faire une rotation d'angle $b$, puis une rotation d'angle $a$, revient à une rotation d'angle $a+b$. Les matrices satisfont $R_aR_b=R_{a+b}$.
\begin{demonstration}
En multipliant les deux matrices, nous obtenons
\[
\begin{aligned}
R_aR_b
&=\begin{pmatrix}\cos a&-\sin a\\\sin a&\cos a\end{pmatrix}
\begin{pmatrix}\cos b&-\sin b\\\sin b&\cos b\end{pmatrix}\\
&=\begin{pmatrix}
\cos a\cos b-\sin a\sin b&-(\sin a\cos b+\cos a\sin b)\\
\sin a\cos b+\cos a\sin b&\cos a\cos b-\sin a\sin b
\end{pmatrix}.
\end{aligned}
\]
Mais cette matrice est $R_{a+b}$, dont les coefficients de la première colonne sont $(\cos(a+b),\sin(a+b))$. L'identification de ces coefficients démontre les deux formules :
\[
\boxed{\begin{aligned}
\cos(a+b)&=\cos a\cos b-\sin a\sin b,\\
\sin(a+b)&=\sin a\cos b+\cos a\sin b.
\end{aligned}}
\]
La preuve est valable pour tous réels $a,b$, même négatifs. $\square$
\end{demonstration}
\section{Formules de soustraction}
En remplaçant $b$ par $-b$ dans les formules d'addition et en utilisant $\cos(-b)=\cos b$, $\sin(-b)=-\sin b$ :
\[
\boxed{\begin{aligned}
\cos(a-b)&=\cos a\cos b+\sin a\sin b,\\
\sin(a-b)&=\sin a\cos b-\cos a\sin b.
\end{aligned}}
\]
\textbf{Application : calculer $\sin(15^\circ)$ exactement.} Comme $15^\circ=60^\circ-45^\circ$ :
\[
\sin\frac{\pi}{12}
=\sin\frac\pi3\cos\frac\pi4-\cos\frac\pi3\sin\frac\pi4
=\boxed{\frac{\sqrt6-\sqrt2}{4}}.
\]
De même $\cos(\pi/12)=(\sqrt6+\sqrt2)/4$.
\end{document}
